%----------------------------------------------------------------------------------------
% METADATA BLOCK (YAML)  — see inbox/README.md
%----------------------------------------------------------------------------------------
% uid: ""
% school: "lycee"
% level: "3eme"
% section: "Economie et Gestion"
% chapter: "Statistiques"
% title: "Moyennes d'une classe"
% difficulty: "Facile"
% tags: ["statistiques", "moyenne", "moyenne-ponderee"]
% provenance: "own"
% batch: "2026-11-19-3eco-statistiques-n5"
%----------------------------------------------------------------------------------------

\newcommand{\R}{\mathbb{R}}
\newcommand{\vect}[1]{\overrightarrow{#1}}

% ---STATEMENT---
Dans une classe de $24$ élèves, la moyenne d'un devoir est $11{,}5$. Les $12$ garçons ont une moyenne de $10{,}5$.
\begin{enumerate}
  \item Calculer la moyenne des filles.
  \item Un élève absent passe le devoir plus tard et obtient $14$. Calculer la nouvelle moyenne de la classe.
  \item Quelle note aurait-il fallu à cet élève pour que la moyenne de la classe devienne $12$ ? Est-ce possible ?
\end{enumerate}

% ---SOLUTION---
\begin{prop}[Moyenne et total]
\ib{Total des notes $=$ effectif $\times$ moyenne}. La moyenne de deux groupes réunis se calcule avec les totaux, pas en faisant la moyenne des deux moyennes.
\end{prop}
\begin{enumerate}[itemsep=14pt]
  \item Total de la classe : $24\times11{,}5=276$. \quad Total des garçons : $12\times10{,}5=126$.

  Total des $12$ filles : $276-126=150$, donc
  \[\overline x_{\text{filles}}=\dfrac{150}{12}=12{,}5\]

  \item Il y a maintenant $25$ élèves et le total devient $276+14=290$ :
  \[\overline x=\dfrac{290}{25}=11{,}6\]

  \item Pour une moyenne de $12$ avec $25$ élèves, il faut un total de $25\times12=300$.

  La note nécessaire est $300-276=24$.

  Une note sur $20$ ne peut pas dépasser $20$ : \textbf{c'est impossible}.
\end{enumerate}
